\documentclass[a4paper,12pt]{article}
\usepackage[latin1]{inputenc}
\usepackage[spanish]{babel}

\usepackage{graphics,graphicx,color}

\usepackage{setspace} % paquete para interlineado

\usepackage{geometry}
\geometry{verbose,a4paper,tmargin=20mm,bmargin=20mm,lmargin=20mm,rmargin=20mm}

\def\membrete{
\vspace*{-10mm}
\hspace{90mm}
\vbox{\hsize = 70mm
\begin{center}\footnotesize\sf
\includegraphics[scale=0.2]{LogoASI-Granada}\\
{\bf Asociaci\'on de J\'ovenes\\ Investigadores de Granada}\\
Facultad de Ciencias\\
Avda. Fuentenueva, s/n\\
18071 Granada\\
Granada
\end{center}}}


\begin{document}
\begin{onehalfspace}

\membrete


\begin{center}
\Large{\bf Estatutos de la\\ Asociaci\'on de J\'ovenes Investigadores de Granada}
\end{center}

\textit{ }

\begin{center}
\textbf{CAP\'ITULO I. Disposiciones Generales}
\end{center}

\textit{} \\
\textbf{Art\'iculo 1.-} Con la denominaci\'on de ASI, Asociaci\'on de J\'ovenes Investigadores de Granada, se constituy\'o en la localidad de Granada, en la provincia de Granada, una entidad de naturaleza asociativa y sin \'animo de lucro, al amparo de la Ley 191/1964, de 24 de diciembre, reguladora de las Asociaciones. En cumplimiento de la Ley Org\'anica 1/2002, de 22 de marzo, reguladora del Derecho de Asociaci\'on, y de su Disposici\'on Transitoria Primera, la asociaci\'on pasa a regirse por los presentes Estatutos, adaptados a dicha Ley.\\
\textit{ }\\
\textit{ }\\
\textbf{Art\'iculo 2.-} La Asociaci\'on constituida tiene personalidad jur\'idica propia y plena capacidad de obrar, pudiendo realizar, en consecuencia, todos aquellos actos que sean necesarios para el cumplimiento de la finalidad para la que ha sido creada, con sujeci\'on a lo establecido en el ordenamiento jur\'idico.\\
\textit{ }\\
\textit{ }\\
\textbf{Art\'iculo 3.-} La Asociaci\'on que se crea tiene nacionalidad espa\~nola. El domicilio social de la Asociaci\'on radicar\'a en la Facultad de Ciencias de la Universidad de Granada, Avda. Fuentenueva, s/n, de Granada, provincia de Granada, C.P. 18071. El cambio de domicilio requerir\'a acuerdo de la Asamblea General, convocada espec\'ificamente con tal objeto, y la modificaci\'on de los presentes Estatutos.\\
\textit{ }\\
\textit{ }\\
\textbf{Art\'iculo 4.-} El \'ambito territorial de acci\'on de la Asociaci\'on se corresponde con el territorio de la provincia de Granada y de las Ciudades Aut\'onomas de Ceuta y Melilla.\\
\textit{ }\\
\textit{ }\\
\textbf{Art\'iculo 5.-} La Asociaci\'on se constituye por tiempo indefinido, y podr\'an nuevos socios causar alta en ella, o baja los antiguos, sin necesidad de nueva constituci\'on.

\textit{ }\\

\begin{center}
\textbf{CAP\'ITULO II. Objeto de la Asociaci\'on}
\end{center}

\textit{ }\\
\textbf{Art\'iculo 6.-} Los fines de la Asociaci\'on ser\'an los siguientes:
\begin{itemize}
\item [a)] Defender los derechos del Personal de investigaci�n.
\item [b)] Promover todo tipo de acciones l\'icitas encaminadas a lograr una mejora continua de las condiciones en las que el Personal Investigador en fase inicial y perfeccionamiento desarrolla su tarea.
\item [c)] Promover todo tipo de acciones l\'icitas encaminadas a lograr para el Personal Investigador en fase inicial y perfeccionamiento el reconocimiento de los derechos inherentes a todo trabajador, seg\'un lo establecido en el Estatuto de  los Trabajadores y normas complementarias.
\item [d)] Fomentar el desarrollo y difusi\'on de la ciencia, la cultura y la investigaci\'on cient\'ifica en Espa\~na.
\end{itemize}

\textit{ }\\
\begin{center}
\textbf{CAP\'ITULO III. De los \'organos de gobierno y forma de administraci\'on}
\end{center}
\textit{ }

\textit{SECCI�N 1�. La asamblea General de socios}

\textit{ }\\
\textbf{Art\'iculo 7.-} El \'organo supremo y soberano de la Asociaci\'on es la Asamblea General de Socios o Asamblea General, integrada por la totalidad de los socios, que se hallen en uso pleno de sus derechos sociales. Adopta sus acuerdos por el principio mayoritario o de democracia interna. Las Asambleas podr\'an tener car\'acter ordinario y extraordinario, en la forma y competencias que se indican en los presentes Estatutos. La Asamblea General deber\'a reunirse, al menos, dos veces al a\~no con car\'acter ordinario.\\
\textit{ }\\
\textit{ }\\
\textbf{Art\'iculo 8.-} Las Asambleas ser\'an convocadas por el Presidente de la Asociaci\'on, por iniciativa propia, por acuerdo de la Junta Directiva o por solicitud firmada por al menos el 20\% del n\'umero legal de socios.\\
Acordado por la Junta Directiva la convocatoria de una Asamblea General, el Presidente habr\'a de convocarla en el plazo m\'aximo de quince d\'ias naturales, para su celebraci\'on. La fecha de celebraci\'on de la asamblea no podr\'a ser superior a un mes a contar desde la fecha del acuerdo.\\
La solicitud de convocatoria efectuada por los socios habr\'a de contener expresamente el orden de d\'ia de la sesi\'on, adjuntando los documentos o informaci\'on que fuere necesaria para la adopci\'on de los acuerdos, si dicha documentaci\'on o informaci\'on hubiere de ser tenida en cuenta para ello.\\
La solicitud habr\'a de ser presentada ante el Secretario de la Asociaci\'on, quien sellar\'a una copia para su entrega al presentador de aqu\'ella.\\
El Secretario de la Asociaci\'on, despu\'es de comprobar los requisitos formales (n\'umero de socios, convocatoria y documentaci\'on, en su caso), dar\'a cuenta inmediata al presidente, para que, en el plazo de quince d\'ias desde su presentaci\'on, convoque la Asamblea que habr\'a de celebrarse dentro del mes siguiente a la fecha de la presentaci\'on. Si la solicitud adoleciere de los requisitos formales antecitados, el Secretario tendr\'a por no formulada la solicitud, procediendo a su archivo con comunicaci\'on al socio que encabece la lista o firmas.\\
Si el Presidente no convocare en el plazo de los quince d\'ias subsiguientes o convocare la Asamblea dentro del plazo para su celebraci\'on con posterioridad al mes desde la solicitud, los promotores estar\'an legitimados para proceder a la convocatoria de la Asamblea General, expresando dichos extremos en la convocatoria, que ir\'a firmada por la persona que encabece las firmas o lista de la solicitud.\\
\textit{ }\\
\textit{ }\\
\textbf{Art\'iculo 9.-} La convocatoria efectuada por las personas legitimadas para ello, de conformidad con lo dispuesto en el art\'iculo anterior, habr\'a de ser comunicada con una antelaci\'on de al menos quince d\'ias a la celebraci\'on de la Asamblea, y ser\'a comunicada a los socios en forma acordada por la Junta Directiva.\\
La convocatoria deber\'a contener el orden del d\'ia, as\'i como el lugar, fecha y hora de su celebraci\'on.\\
La documentaci\'on necesaria e informaci\'on que haya de ser tenida en cuenta para la adopci\'on de los acuerdos, habr\'a de estar a disposici\'on de los socios en la Secretar\'ia de la Asociaci\'on, con una antelaci\'on m\'inima de quince d\'ias a la celebraci\'on de la Asamblea, la cual podr\'a ser examinada por aqu\'ellos en la expresada Secretar\'ia.\\
\textit{ }\\
\textit{ }\\
\textit{ }\\
\textit{ }\\
\textbf{Art\'iculo 10.-} La Asamblea General Ordinaria habr\'a de convocarse durante los tres primeros meses de cada a\~no, al objeto de tratar los siguientes puntos del orden del d\'ia:
\begin{itemize}
\item [1.-] Lectura y aprobaci\'on, si procede, del acta de la sesi\'on anterior (sea Asamblea General Ordinaria o Extraordinaria).
\item [2.-] Examen y aprobaci\'on, si procediere, de las Cuentas del ejercicio anterior.
\item [3.-] Examen de la memoria de actividades y aprobaci\'on, si procediere, de la gesti\'on de la Junta Directiva.
\item [4.-] Aprobaci\'on, si procediere, del Programa de Actividades.
\end{itemize}
\textit{ }\\
\textit{ }\\
\textbf{Art\'iculo 11.-} Se requerir\'a la convocatoria de Asamblea General Extraordinaria, para tratar de los siguientes aspectos:
\begin{itemize}
\item [1.-] Modificaci\'on parcial o total de los Estatutos.
\item [2.-] Disoluci\'on de la Asociaci\'on.
\item [3.-] Nombramiento de la Junta Directiva.
\item [4.-] Disposici\'on y Enajenaci\'on de Bienes.
\item [5.-] Constituci\'on de una Federaci\'on, Confederaci\'on o Uni\'on de Asociaciones o su integraci\'on en ella si ya existiere.
\end{itemize}
\textit{ }\\
\textit{ }\\
\textbf{Art\'iculo 12.-} Las Asambleas, tanto ordinarias como extraordinarias, quedar\'an v\'alidamente constituidas, previa convocatoria efectuada con una antelaci\'on de quince d\'ias, cuando concurran a ellas, presentes o representados, un tercio de los asociados con derecho a voto; en segunda convocatoria, cualquiera que sea el n\'umero de asociados con derecho a voto.\\
Para el c\'omputo de socios o n\'umero de votos total, las representaciones habr\'an de presentarse al Secretario con inmediaci\'on al inicio de la sesi\'on.\\
Actuar\'an como presidente de la Asamblea, el presidente de la asociaci\'on, y como secretario de la Asamblea, el secretario de la asociaci\'on.\\
\textit{ }\\
\textit{ }\\
\textit{ }\\
\textbf{Art\'iculo 13.-} Todos los asuntos se debatir\'an y votar\'an conforme figuren en el orden del d\'ia. El presidente iniciar\'a el debate, abriendo un primer turno de intervenciones en el que se har\'a uso de la palabra previa su autorizaci\'on. El presidente moderar\'a los debates, pudiendo abrir un segundo turno de intervenciones o conceder las palabras por alusiones.\\
Los acuerdos de la Asamblea General se adoptar\'an por mayor\'ia simple de las personas presentes o representadas, cuando los votos afirmativos superen a los negativos.\\
No obstante, requerir\'an mayor\'ia cualificada de las personas presentes o representadas, que resultar\'a cuando los votos afirmativos superen la mitad, los acuerdos relativos a disoluci\'on de la asociaci\'on, modificaci\'on de los Estatutos, disposici\'on o enajenaci\'on de bienes y remuneraci\'on de los miembros de la Junta Directiva.\\
Los acuerdos de la Asamblea General que afecten a la denominaci\'on de la Asociaci\'on, domicilio, fines y actividades estatutarias, \'ambito de actuaci\'on, designaci\'on de los miembros de la Junta Directiva, apertura y cierre de delegaciones, constituci\'on de federaciones, confederaciones y uniones, disoluci\'on, o los de modificaciones estatutarias, se comunicar\'an al Registro de Asociaciones para su inscripci\'on, en el plazo de un mes desde que se produzca el acuerdo.\\
\textit{ }\\
\textit{ }\\
\textbf{Art\'iculo 14.-} La representaci\'on o delegaci\'on de voto s\'olo ser\'a v\'alida para la sesi\'on o convocatoria por la que se expida, siendo nula cualquier delegaci\'on o representaci\'on indefinida. Habr\'a de hacerse constar por escrito, con indicaci\'on de los datos personales y n\'umero de DNI del delegante y representado, y firmado y rubricado por ambos.

\textit{ }
\textit{ }\\

\textit{SECCI�N 2�. De los \'organos de representaci\'on}

\textit{ }\\
\textbf{Art\'iculo 15.-} La Junta Directiva es el \'organo colegiado de gobierno, representaci\'on y administraci\'on de la Asociaci\'on, sin perjuicio de las potestades de la Asamblea General como \'organo soberano. Su mandato ser\'a de un a\~no, pudiendo sus miembros ser reelegidos una sola vez.\\
\textit{ }\\
\textit{ }\\
\textbf{Art\'iculo 16.-} La Junta Directiva estar\'a formada por un Presidente, un Vicepresidente, un Secretario, un Tesorero y un n\'umero indefinido de vocales, que ser\'a fijado por la Asamblea General, y que en cualquier caso no ser\'a inferior a dos. Todos los miembros de la Junta Directiva ser\'an designados y revocados por la Asamblea General. El ejercicio del cargo ser\'a personal, sin que pueda delegarse el voto en las sesiones de la Junta Directiva.\\
\textit{ }\\
\textit{ }\\
\textit{ }\\
\textit{ }\\
\textbf{Art\'iculo 17.-} Para ser miembro de la Junta Directiva ser\'an requisitos imprescindibles ser mayor de edad, estar en pleno uso de los derechos civiles y no estar incurso en los motivos de incompatibilidad establecidos en la legislaci\'on vigente.\\
Los miembros de la Junta Directiva ser\'an elegidos, entre los socios, en Asamblea General Extraordinaria, de conformidad con lo dispuesto en el art\'iculo 11.\\
Una vez convocada la Asamblea General para la designaci\'on de los miembros de la Junta Directiva, los socios que pretendan ejercer su derecho de elegibilidad, habr\'an de presentar su candidatura con una antelaci\'on, como m\'inimo, de veinticuatro horas a la celebraci\'on de la Asamblea. Las candidaturas se presentar\'an en forma de lista cerrada, indicando las personas que integran la candidatura, as\'i como el cargo al que aspiran dentro de la Junta Directiva.\\
Al comienzo de la Asamblea General para la designaci\'on de los miembros de la Junta Directiva, la anterior Junta Directiva, constituida en Comisi\'on Electoral, dar\'a a conocer las candidaturas, aceptando y resolviendo las reclamaciones que sobre las mismas se pudieran derivar. Producida una vacante, provisionalmente, la Junta Directiva podr\'a designar a otro miembro de la misma para su sustituci\'on, hasta que se produzca la elecci\'on del vocal correspondiente por la Asamblea General en la primera sesi\'on que se convoque.\\
\textit{ }\\
\textit{ }\\
\textbf{Art\'iculo 18.-} Los miembros de la Junta Directiva cesar\'an en sus respectivos cargos por las siguientes causas:
\begin{itemize}
\item [a)] Por muerte o declaraci\'on de fallecimiento.
\item [b)] Por incapacidad, inhabilitaci\'on o incompatibilidad, de acuerdo con lo establecido en el ordenamiento jur\'idico.
\item [c)] Por resoluci\'on judicial.
\item [d)] Por transcurso del periodo de su mandato. No obstante, hasta tanto no se proceda a la Asamblea General subsiguiente para la elecci\'on de la nueva Junta Directiva, aqu\'ella continuar\'a en funciones, debi\'endose expresar dicho car\'acter en cuantos documentos hubieren de firmar en funci\'on a los respectivos cargos.
\item [e)] Por renuncia.
\item [f)] Por acuerdo adoptado con las formalidades estatutarias, en cualquier momento, por la Asamblea General.
\item [g)] Por la p\'erdida de la condici\'on de socio.
\end{itemize}
\textit{ }\\
Los ceses y nombramientos habr\'an de ser comunicados al Registro de Asociaciones, para su debida constancia y publicidad.\\
\textit{ }\\
\textit{ }\\
\textbf{Art\'iculo 19.-} Corresponde al Presidente:
\begin{itemize}
\item [a)] Ostentar la representaci\'on de la Asociaci\'on ante toda clase de personas, autoridades y entidades p\'ublicas o privadas.
\item [b)] Convocar las reuniones de la Junta Directiva y de la Asamblea General, presidirlas, dirigir sus debates, suspender y levantar las sesiones.
\item [c)] Ejecutar los acuerdos de la Junta Directiva y de la Asamblea General, pudiendo para ello realizar toda clase de actos y contratos y firmar aquellos documentos necesarios a tal fin; sin perjuicio de que por cada \'organo en el ejercicio de sus competencias, al adoptar los acuerdos se faculte expresamente para su ejecuci\'on a cualquier otro miembro de la Junta Directiva.
\item [d)] Cumplir y hacer cumplir los acuerdos de la Junta Directiva y Asamblea General.
\item [e)] Ordenar los gastos y pagos de la Asociaci\'on.
\item [f)] Dirimir con su voto los empates.
\item [g)] Visar las actas y certificaciones de los acuerdos de la Junta Directiva y Asamblea General.
\item [h)] Ejercer cuantas otras funciones sean inherentes a su condici\'on de Presidente de la Junta Directiva y de la Asociaci\'on.
\end{itemize}
\textit{ }\\
\textit{ }\\
\textbf{Art\'iculo 20.-} Corresponde al Vicepresidente:
\begin{itemize}
\item [a)]  Sustituir al Presidente en los casos de vacante, ausencia o enfermedad.
\item [b)]  Colaborar con el Presidente y ejercer las atribuciones que le sean delegadas o encomendadas por \'este.
\textit{ }\\
\textit{ }\\
\end{itemize}
\textbf{Art\'iculo 21.-} Corresponde al Secretario:
\begin{itemize}
\item [a)] Asistir a las sesiones de la Junta Directiva y Asamblea y redactar y autorizar las actas de aqu\'ellas.
\item [b)] Efectuar la convocatoria de las sesiones de la Junta Directiva y Asamblea, por orden del Presidente, as\'i como las citaciones de los miembros de aqu\'ella y socios de \'esta.
\item [c)] Dar cuenta inmediata al Presidente de la solicitud de convocatoria efectuada por los socios en la forma prevista en el art\'iculo 8 de los presentes Estatutos.
\item [d)] Recibir los actos de comunicaci\'on de los miembros de la Junta Directiva con relaci\'on a \'esta y de los socios y, por tanto, las notificaciones, peticiones de datos, rectificaciones, certificaciones o cualquiera otra clase de escritos de los que deba tener conocimiento.
\item [e)] Preparar el despacho de los asuntos, y por lo tanto con la documentaci\'on correspondiente que hubiere de ser utilizada o tenida en cuenta.
\item [f)] Expedir certificaciones de los acuerdos aprobados y cualesquiera otras certificaciones, con el visto bueno del Presidente, as\'i como los informes que fueren necesarios.
\item [g)] Tener bajo su responsabilidad y custodia el Archivo, documentos y Libros de la Asociaci\'on, a excepci\'on del/los libros de contabilidad.
\item [h)] Cualesquiera otras funciones inherentes a su condici\'on de Secretario.
\end{itemize}
\textit{ }\\
En los casos de ausencia o enfermedad y, en general, cuando concurra alguna causa justificada, el Secretario ser\'a sustituido por el vocal de menor edad.\\
\textit{ }\\
\textit{ }\\
\textbf{Art\'iculo 22.-} Corresponde al Tesorero:
\begin{itemize}
\item [a)] Recaudar los fondos de la Asociaci\'on, custodiarlos e invertirlos en la forma determinada por la Junta Directiva.
\item [b)] Efectuar los pagos, con el visto bueno del Presidente.
\item [c)] Intervenir con su firma todos los documentos de cobros y pagos, con el conforme del Presidente.
\item [d)] La llevanza de los libros de contabilidad y el cumplimiento de las obligaciones fiscales, en plazo y forma, de la Asociaci\'on.
\item [e)] La elaboraci\'on del Estado General de Cuentas para su aprobaci\'on anual por la Asamblea.
\item [f)] Cualesquiera otras inherentes a su condici\'on de tesorero, como responsable de la gesti\'on econ\'omica financiera.
\end{itemize}
\textit{ }\\
\textit{ }\\
\textbf{Art\'iculo 23.-} De conformidad con el art\'iculo 16, en la Junta Directiva habr\'a al menos dos vocales:
\begin{itemize}
\item  Un vocal que actuar\'a como portavoz y representante de la Asociaci\'on ante la Federaci\'on de J\'ovenes Investigadores.
\item  Un vocal que actuar\'a como portavoz de la asociaci\'on ante los medios de comunicaci\'on.
\end{itemize}
\textit{ }\\
La Asamblea General podr\'a designar adem\'as nuevos vocales, delimitando en el momento de su designaci\'on las funciones espec\'ificas que tendr\'an.\\
\textit{ }\\
\textit{ }\\
\textit{ }\\
\textit{ }\\
\textbf{Art\'iculo 24.-} De forma gen\'erica, y sin menoscabo de sus funciones espec\'ificas, corresponde a los vocales:
\begin{itemize}
\item [a)] Recibir la convocatoria de la sesi\'on de Junta Directiva con la antelaci\'on fijada en los presentes Estatutos, conteniendo aqu\'ella el orden del d\'ia.
\item [b)] Participar en el debate de las sesiones.
\item [c)] Ejercer su derecho al voto y formular su voto particular, as\'i como expresar el sentido de su voto y los motivos que lo justifican.
\item [d)] Formular ruegos y preguntas.
\item [e)] Obtener la informaci\'on precisa para el cumplimiento de las funciones que le fueren asignadas.
\end{itemize}
\textit{ }\\
\textit{ }\\
\textbf{Art\'iculo 25.-} Convocatorias y sesiones de la Junta Directiva
\begin{itemize}
\item [1.-] Para la v\'alida constituci\'on de la Junta Directiva, a efectos de la celebraci\'on de sesiones, deliberaciones y adopci\'on de acuerdos, deber\'an estar presentes la mitad de sus miembros requiri\'endose, necesariamente, la presencia del Presidente y del Secretario o de quienes les sustituyan.
\item [2.-] La Junta Directiva se reunir\'a, al menos, una vez al trimestre y cuantas veces como sea preciso para la buena marcha de la Asociaci\'on, por convocatoria realizada por el Presidente, a iniciativa propia o de cualesquiera de sus miembros.
\item [3.-] La convocatoria , con sus elementos formales (orden del d\'ia, lugar y fecha...), se har\'a llegar con una antelaci\'on m\'inima de 48 horas a su celebraci\'on.
\item [4.-] Las deliberaciones seguir\'an el mismo r\'egimen se\~nalado en el art\'iculo 13 para la Asamblea General. Los acuerdos se adoptar\'an por mayor\'ia simple de los votos emitidos, dirimiendo el voto del Presidente en caso de empate.
\item [5.-] No podr\'a adoptarse acuerdo alguno que no figure en el orden del d\'ia, salvo que, estando presentes la totalidad de los miembros que componen la Junta Directiva, lo acuerden por unanimidad.
\item [6.-] Igualmente quedar\'a v\'alidamente constituida la Junta Directiva sin convocatoria previa, cuando estando presentes todos y cada uno de los miembros, as\'i se acuerde por unanimidad, respetando lo mencionado en el apartado anterior en cuanto a los acuerdos. Las Juntas as\'i constituidas recibir\'an la denominaci\'on de Junta Directiva Universal.
\item [7.-] A las sesiones de la Junta Directiva podr\'an asistir aquellas personas con funciones de asesoramiento, previamente citadas o invitadas por el Presidente, con voz y sin voto para mejor acierto en sus deliberaciones.
\end{itemize}
\textit{ }\\
\textit{ }\\
\textbf{Art\'iculo 26.-} La Junta Directiva tendr\'a las siguientes atribuciones:
\begin{itemize}
\item [a)] Confeccionar el Plan de Actividades.
\item [b)] Otorgar apoderamientos generales o especiales.
\item [c)] Organizar y desarrollar las actividades aprobadas por la Asamblea General.
\item [d)] Aprobar el estado de Cuentas elaborado por el Tesorero, para su aprobaci\'on definitiva, si procediere, por la Asamblea General.
\item [e)] Elaborar la Memoria anual de actividades para su informe a la Asamblea General.
\item [f)] Creaci\'on de Comisiones de Trabajo que estime conveniente, para el desarrollo de las funciones encomendadas y las actividades aprobadas, as\'i como para cualesquiera otras cuestiones derivadas del cumplimiento de los fines sociales. Dichas comisiones regular\'an su funcionamiento interno en la forma que se acuerden por \'estas en su primera sesi\'on constitutiva.
\item [g)] Resolver las solicitudes relativas a la admisi\'on de socios.
\end{itemize}
\textit{ }\\
\textit{ }\\
\textbf{Art\'iculo 27.-} Son obligaciones de los miembros de la Junta Directiva, a t\'itulo enunciativo, cumplir y hacer cumplir los fines de la Asociaci\'on, concurrir a las reuniones a las que sean convocados, desempe\~nar el cargo con la debida diligencia de un representante leal y cumplir en sus actuaciones con lo determinado en las disposiciones legales vigentes y en los presentes Estatutos.\\
Los miembros de la Junta Directiva responder\'an frente a la Asociaci\'on de los da\~nos y perjuicios que causen por actos contrarios a la Ley o a los Estatutos o por los realizados negligentemente. Quedar\'an exentos de responsabilidad quienes se opusieren expresamente al acuerdo determinante de tales actos o no hubiere participado en su adopci\'on.\\
\textit{ }\\
\textit{ }\\
\textbf{Art\'iculo 28.-} Los miembros de la Junta Directiva ejercer\'an su cargo gratuitamente, sin que en ning\'un caso puedan recibir retribuci\'on por el desempe\~no de su funci\'on, sin perjuicio del derecho a ser reembolsados en los gastos ocasionados en ejercicio de los cargos, siempre que \'estos se encuentren debida y formalmente justificados.

\textit{ }

\textit{SECCI�N 3�. Disposiciones comunes a los \'organos}

\textit{ }\\
\textbf{Art\'iculo 29.-} De las actas.
\begin{itemize}
\item [1.-] De cada sesi\'on que celebren la Asamblea General y Junta Directiva   se   levantar\'a    acta   por   el   Secretario,  que  especificar\'a necesariamente el qu\'orum necesario para la v\'alida constituci\'on (en el caso de la Junta Directiva se especificar\'a necesariamente los asistentes), el orden del d\'ia de la reuni\'on, las circunstancias del lugar y tiempo en que se han celebrado, los puntos principales de las deliberaciones, as\'i como el contenido de los acuerdos adoptados.
\item [2.-] En el acta figurar\'a, a solicitud de los respectivos miembros y/o socios, el voto contrario al acuerdo adoptado, su abstenci\'on y los motivos que la justifiquen o el sentido de su voto favorable. Asimismo, cualquier miembro tiene derecho a solicitar la transcripci\'on \'integra de su intervenci\'on o propuesta, siempre que aporte en el acto o en el plazo de cuarenta y ocho horas el texto que se corresponda fielmente con su intervenci\'on, haci\'endose as\'i constar en el acta o uni\'endose copia a la misma.
\item [3.-] Las actas se aprobar\'an en la misma o en la siguiente sesi\'on, pudiendo no obstante emitir el Secretario certificaci\'on sobre los acuerdos espec\'ificos que se hayan adoptado, sin perjuicio de la ulterior aprobaci\'on del acta. En las certificaciones de acuerdos adoptados emitidas con anterioridad a la aprobaci\'on del acta se har\'a constar expresamente tal circunstancia.
\item [4.-] Las Actas ser\'an firmadas por el Secretario y visadas por el Presidente.
\end{itemize}
\textit{ }\\
\textit{ }\\
\textbf{Art\'iculo 30.-} Los acuerdos de la Asamblea General y Junta Directiva podr\'an ser impugnados ante el orden jurisdiccional civil en la forma legalmente establecida. Los asociados podr\'an impugnar los acuerdos y actuaciones de la asociaci\'on que estimen contrarios a los Estatutos dentro del plazo de cuarenta d\'ias, a partir de la fecha de adopci\'on de los mismos, instando su rectificaci\'on o anulaci\'on y la suspensi\'on preventiva en su caso, o acumulando ambas pretensiones por los tr\'amites establecidos en la Ley de Enjuiciamiento Civil. En tanto se resuelven las contiendas de orden interno que puedan suscitarse en la asociaci\'on, las solicitudes de constancia registral que se formulen sobre las cuestiones controvertidas s\'olo dar\'an lugar a anotaciones provisionales.

\textit{ }\\
\begin{center}
\textbf{CAP\'ITULO IV. Procedimiento de admisi\'on y p\'erdida de la cualidad de socio}
\end{center}
\textit{ }\\
\textbf{Art\'iculo 31.-} Para adquirir la condici\'on de socio se requiere ser persona f\'isica o jur\'idica, y estar interesada en los fines de la Asociaci\'on.
\textit{ }\\
La solicitud para adquirir la condici\'on de socio debe aceptarse por la Junta Directiva.\\
\textit{ }\\
\textit{ }\\
\textbf{Art\'iculo 32.-} La condici\'on de socio se perder\'a por alguna de las siguientes causas:
\begin{itemize}
\item [a)] Por renuncia voluntaria, comunicada por escrito a la Junta Directiva.
\item [b)] Por impago reiterado de la cuota de la asociaci\'on.
\item [c)] Por incumplimiento grave de los presentes Estatutos o de los acuerdos v\'alidamente adoptados por los \'organos sociales.
\end{itemize}
\textit{ }\\

\begin{center}
\textbf{CAP\'ITULO V. Derechos y Deberes de los Socios}
\end{center}
\textbf{Art\'iculo 33.-} Son derechos de los socios:
\begin{itemize}
\item [a)] Participar en las actividades de la asociaci\'on y en los \'organos de gobierno y representaci\'on, a ejercer el derecho de voto, as\'i como a asistir a la Asamblea General, de acuerdo con los Estatutos.
\item [b)] Ser informado acerca de la composici\'on de los \'organos de gobierno y representaci\'on de la asociaci\'on, de su estado de cuentas y del desarrollo de su actividad.
\item [c)] Ser o\'ido con car\'acter previo a la adopci\'on de medidas disciplinarias contra \'el y ser informado de los hechos que den lugar a tales medidas, debiendo ser motivado el acuerdo que, en su caso, imponga la sanci\'on.
\item [d)] Acceder a la documentaci\'on de la asociaci\'on, a trav\'es de la Junta Directiva.
\item [e)] Usar los bienes e instalaciones de uso com\'un de la Asociaci\'on, con respeto a igual derecho del resto de los socios.
\end{itemize}
\textit{ }\\
\textit{ }\\
\textbf{Art\'iculo 34.-} Son deberes de los socios:
\begin{itemize}
\item [a)] Compartir las finalidades de la asociaci\'on y colaborar para la consecuci\'on de las mismas.
\item [b)] Pagar las cuotas, derramas y otras aportaciones que, con arreglo a los Estatutos, puedan corresponder a cada socio.
\item [c)] Cumplir el resto de obligaciones que resulten de las disposiciones estatutarias.
\item [d)] Acatar y cumplir los acuerdos v\'alidamente adoptados por la Junta Directiva y la Asamblea General.
\end{itemize}
\textit{ }\\
Sin perjuicio de la p\'erdida de la condici\'on de socio por impago de las cuotas sociales, \'interin se procede a su expulsi\'on, el socio tendr\'a en suspenso el derecho de sufragio activo y pasivo. Dicha suspensi\'on del derecho se producir\'a con el impago de una sola de las cuotas y mientras se proceda a su regularizaci\'on o a la p\'erdida definitiva de la condici\'on de socio.

\textit{ }\\
\begin{center}
\textbf{CAP\'ITULO VI. R\'egimen econ\'omico}
\end{center}
\textit{ }\\
\textbf{Art\'iculo 35.-} El patrimonio fundacional de la Asociaci\'on en el momento de su constituci\'on es de cero euros.\\
\textit{ }\\
\textit{ }\\
\textbf{Art\'iculo 36.-} La Asociaci\'on deber\'a figurar como titular de todos los bienes y derechos que integran su patrimonio, los cuales se har\'an constar en su inventario y se inscribir\'an, en su caso, en los Registros P\'ublicos correspondientes.\\
\textit{ }\\
\textit{ }\\
\textbf{Art\'iculo 37.-} La Asociaci\'on, para el desarrollo de sus actividades, se financiar\'a con:
\begin{itemize}
\item [a)] Los recursos que provengan del rendimiento de su patrimonio, en su caso.
\item [b)] Las cuotas de los socios, ordinarias o extraordinarias.
\item [c)] Los donativos o subvenciones que pudieran ser concedidas por personas f\'isicas o jur\'idicas, p\'ublicas o privadas.
\item [d)] Donaciones, herencias o legados, aceptadas por la Junta Directiva.
\item [e)] Los ingresos provenientes de sus actividades.
\end{itemize}
\textit{ }\\
Los beneficios obtenidos por la Asociaci\'on, derivados del ejercicio de actividades econ\'omicas, incluidas las prestaciones de servicios, deber\'an destinarse exclusivamente al cumplimiento de sus fines, sin que quepa en ning\'un caso su reparto entre los asociados ni entre sus c\'onyuges o personas que convivan con aqu\'ellos con an\'aloga relaci\'on de afectividad, ni entre sus parientes, ni su cesi\'on gratuita a personas f\'isicas o jur\'idicas con inter\'es lucrativo.\\
\textit{ }\\
\textit{ }\\
\textbf{Art\'iculo 38.-} El ejercicio econ\'omico coincidir\'a con el a\~no natural, por lo que comenzar\'a el 1 de enero y finalizar\'a el 31 de diciembre de cada a\~no.\\
Para la aprobaci\'on de cuotas extraordinarias, habr\'a de convocarse en Asamblea General Extraordinaria, salvo que la Asociaci\'on careciere de liquidez y la disposici\'on y gasto correspondiente fueren urgentes, en cuyo caso bastar\'a la adopci\'on del acuerdo por la Junta Directiva, previo informe del Tesorero y ulterior ratificaci\'on en Asamblea General, que habr\'a de aprobarse en el plazo de treinta d\'ias siguientes a la adopci\'on del acuerdo por la Junta Directiva.\\
La Asamblea General aprobar\'a anualmente las cuentas de la Asociaci\'on, una vez finalizado el ejercicio presupuestario al que correspondan.\\
La Junta Directiva llevar\'a los correspondientes libros de contabilidad, que permitan obtener una imagen fiel del patrimonio, el resultado y la situaci\'on financiera de la Asociaci\'on.

\textit{ }\\
\begin{center}
\textbf{CAP\'ITULO VII. Disoluci\'on y aplicaci\'on del capital social}
\end{center}
\textit{ }\\
\textit{ }\\
\textbf{Art\'iculo 39.-} La Asociaci\'on se disolver\'a por las siguientes causas:
\begin{itemize}
\item [a)] Por acuerdo adoptado por mayor\'ia cualificada en Asamblea General Extraordinaria.
\item [b)] Por las causas que se determinan en el art\'iculo 39 del C\'odigo Civil.
\item [c)] Por sentencia judicial firme.
\end{itemize}
\textit{ }\\
\textit{ }\\
\textbf{Art\'iculo 40.-} Acordada la disoluci\'on de la Asociaci\'on, se abre el periodo de liquidaci\'on, hasta el fin del cual la entidad conservar\'a su personalidad jur\'idica. Los miembros de la Junta Directiva en el momento de la disoluci\'on se convierten en liquidadores, salvo que los designe expresamente la Asamblea General o el juez que, en su caso, acuerde la disoluci\'on.\\
Corresponde a los liquidadores:
\begin{itemize}
\item [a)] Velar por la integridad del patrimonio de la Asociaci\'on.
\item [b)] Concluir las operaciones pendientes y efectuar las nuevas que sean precisas para la liquidaci\'on.
\item [c)] Cobrar los cr\'editos de la Asociaci\'on.
\item [d)] Liquidar el patrimonio y pagar a los acreedores.
\item [e)] Aplicar los bienes sobrantes a los fines previstos por los Estatutos.
\item [f)] Solicitar la cancelaci\'on de los asientos en el Registro.
\end{itemize}
El patrimonio resultante despu\'es de pagadas las deudas y cargas sociales, se destinar\'a a entidades no lucrativas que persigan fines de inter\'es general an\'alogos a los realizados por la misma.\\
En caso de insolvencia de la Asociaci\'on, la Junta Directiva o, en su caso, los liquidadores, han de promover inmediatamente el oportuno procedimiento concursal ante el juez competente.

\end{onehalfspace}
\end{document}
